\chapter{Interactive GSLTs, and the Site of Interaction}
\label{ch:igslt}
%===============================================================================

\section{Naming the site of interaction}

A GSLT is any multi-sorted rewrite theory. Most of what we want to do requires more:
we want to know \emph{where} the theory's programs meet whatever they compute against.

\begin{definition}[Interactive GSLT]
\label{ob:def:igslt}
A GSLT $G = (\Sigma,E,R)$ is \textbf{interactive} when it carries, over and above the
triple:
\begin{itemize}[leftmargin=2em]
  \item a distinguished \emph{interacting sort} $P$ of programs;
  \item a binary \emph{interaction constructor} $K(P,E)$, representing the interaction
    of a program with an environment $E$ --- an operand of the interacting sort, with
    $E = P$ in the process-calculus instances;
  \item at least one base rewrite rule in $R$ whose left-hand side features $K$ --- the
    \emph{interaction rule} --- generating a labelled transition system on $P$ modulo
    $\equiv$ on which bisimulation is the intended equivalence.
\end{itemize}
\end{definition}

\begin{definition}[The category $\catiGSLT$]
$\catiGSLT$ is the category whose objects are interactive GSLTs and whose morphisms are
theory maps --- sort- and operator-preserving translations respecting $\equiv$ and $R$
--- that are bisimulation-preserving on the interacting sort.
\end{definition}

The constructor $K$ and the $K$-headed rule are exactly what a bare GSLT lacks. A GSLT
is any multi-sorted rewrite theory; an interactive GSLT \emph{names a site of
interaction and a rule that fires there}. The motivating instances span calculi with and
without binding: CCS~\cite{milnerCC} with $P$ the processes, $K = {\mid}$ and
interaction by complementary actions; the rho calculus with $K = {\mid}$ and $R \ni
\textsc{Comm}$; the $\lambda$-calculus with $P$ the terms, $K = \mathrm{App}$, $\equiv$
$= \alpha$ and $R \ni \beta$; and the interaction categories of Abramsky, Gay and
Nagarajan~\cite{AbramskyGayNagarajan1996}, where $K$ is composition along a shared
interface.

\section{The interaction cut}
\label{ob:subsec:cut}

Definition~\ref{ob:def:igslt} asks only that \emph{some} rule mention $K$. In every
instance we care about, that rule has a great deal more shape, and naming the shape is
what makes the constructions of Part~III possible. Write the generating rule of $R$ as
\[
  \frac{(C_p', C_e') = \compute(I, J, C_p, C_e)}
       {K\bigl(K_p(I, C_p),\, K_e(J, C_e)\bigr) \longrightarrow K'\bigl(C_p', C_e'\bigr)}
  \tag{$\dagger$}
\]
distinguishing:

\begin{description}[leftmargin=1.5em,style=nextline]
\item[the interaction constructor $K$,] which brings two operands into contact. In rho
  and CCS, $K = {\mid}$; in $\lambda$, $K = \mathrm{App}$.

\item[two introductions $K_p(I,C_p)$ and $K_e(J,C_e)$,] each a constructor separating an
  \emph{interaction surface} ($I$, $J$) --- the part that must match for the cut to fire
  --- from a \emph{continuation}: $C_p$ the program continuation, the rest of the
  computation, and $C_e$ the environment continuation, the rest of the environment.

\item[the contraction $\compute$,] a partial operation that, on a surface match,
  combines the two continuations into residuals $C_p', C_e'$ repackaged by $K'$.
\end{description}

The factoring is Milner's~\cite{Milner1992}. In the polyadic $\pi$-calculus the input
$\mathrm{for}(y \leftarrow x)P$ is the pair $(x, \lambda y.P)$ and the output $x!(u).Q$
is the pair $(x, [u,Q])$: the subjects $x$ are the surfaces that must match, the objects
$\lambda y.P$ and $[u,Q]$ are the continuations, and the contraction is Milner's
\emph{pseudo-application}, returning $P\{u/y\} \mid Q$.

\begin{center}
\begin{tikzpicture}[
  every node/.style={font=\small},
  box/.style={draw,rounded corners=2pt,minimum height=8mm,minimum width=20mm},
  lbl/.style={font=\scriptsize\itshape}
]
  \node[box] (kp) at (-2.6,0) {$K_p(I, C_p)$};
  \node[box] (ke) at ( 2.6,0) {$K_e(J, C_e)$};
  \node (k) at (0,0) {$K$};
  \draw (kp) -- (k); \draw (k) -- (ke);
  \node[lbl] at (-2.6,0.75) {program};
  \node[lbl] at ( 2.6,0.75) {environment};
  \draw[-{Latex[length=2mm]}] (0,-0.5) -- (0,-1.3)
        node[midway,right,font=\scriptsize] {$\compute(I,J,C_p,C_e)$};
  \node[box,minimum width=34mm] (res) at (0,-1.8) {$K'(C_p', C_e')$};
  \draw[dashed] (-1.1,-0.35) -- (1.1,-0.35);
  \node[lbl] at (0,0.75) {surfaces $I$, $J$ must match};
\end{tikzpicture}
\end{center}

We call $(\dagger)$ an \emph{interaction cut}: two introductions meet along matching
surfaces and a contraction combines their continuations. Both sides are structured ---
the eliminand is not a bare term but an introduction carrying its own continuation.

\begin{remark}[Binding is not the point]
Pseudo-application factors into two parts: the substitution $P\{u/y\}$, by which a
datum migrates from the environment continuation into the program continuation through
a binder, and the release $\mid Q$ of the environment's residual. Binding is therefore
\emph{one mode by which information migrates} under the contraction, not an intrinsic
feature of the cut. CCS realises the same shape with no migration at all, and is in this
sense the paradigmatic instance.
\end{remark}

\section{Surfaces: nominal and structural}
\label{ob:subsec:surfaces}

A surface may be carried in either of two ways, and which one a calculus uses is
dictated entirely by the equational theory of $K$.

It may be carried \emph{nominally}, as an explicit $I$ or $J$ that must match --- as the
subject $x$ does in rho and $\pi$, or as complementary labels do in CCS. Or it may be
carried \emph{structurally}, by the rigidity of $K$: when $K$ is free, carrying no
equations at all, position itself is unforgeable and the context $K([\,],-)$ \emph{is}
the surface. A nominally degenerate --- absent --- surface is then not a missing match
condition but one carried by the geometry of $K$ instead.

The $\lambda$-calculus is exactly this case on the environment side. $\mathrm{App}$ is
free, and its rigidity supplies the surface that no explicit $J$ records. Each
application comprises a location together with a function and an argument; the location
is the surface, the function is the program and the argument is the environment --- or
dually. A variable, we note, is not a surface, though it is closely related to one.

The dividing line is \emph{freeness}, not commutativity. An equation on $K$ is a
congruence step: it rewrites adjacency without firing a cut. Associativity re-brackets
adjacency, a unit law inserts and deletes neighbours, idempotence copies them, and
associative--commutativity dissolves position altogether into a freely mixed soup. Each
forges position, and only a free $K$ admits no such move. We return to the consequences
in Section~\ref{ob:subsec:capabilities}.

\section{Discriminating examples: three machines that are not interactive}
\label{ob:subsec:machines}

It is easy to accumulate examples of a definition and hard to see what the definition
excludes. We therefore work three non-examples in detail. They are not exotic: they are
the machines with which most readers first learned what computation is.

\section{Turing machines}

A Turing machine is unproblematically a GSLT. Its sorts are states, tape symbols, and
configurations; its operators build a configuration from a state, a tape, and a head
position; its rewrites are the transition table together with the mechanics of moving
the head. One can write the presentation down and a machine will accept it.

What it is not is \emph{interactive} in the sense of Definition~\ref{ob:def:igslt}. To make
it interactive one must exhibit an interaction constructor $K(P,E)$ and a $K$-headed
rule. The only plausible candidate for the site of interaction is the contact between
the finite control and the tape: that is, after all, where the machine's steps happen.
But the definition asks that $E$ be an operand of the interacting sort, and the tape is
not of the sort of the control. The candidate cut is heterogeneous, and $K(P,E)$ has no
well-sorted instance.

\begin{nonexample}[Turing machines]
\label{ob:nonex:turing}
The Turing machine has no naive presentation as an interactive GSLT. The cut would be
between automaton and tape, which are of different kinds.
\end{nonexample}

The failure is worth dwelling on because it is not a failure of expressive power. A
Turing machine can compute anything a rho calculus term can compute. What it lacks is a
\emph{site}: a place in the syntax where two things of the same kind are brought into
contact, such that one could ask what else might have been brought into contact there
instead. The absence of such a site is precisely the absence of an environment in which
the machine could be placed --- which is another way of saying that a Turing machine is
a model of a closed computation, and interactivity is about open ones.

\section{Moore and Mealy machines}

The transducers fail in the same way and differ from each other in an instructive
second respect.

A Moore machine emits an output determined by its current state alone; a Mealy machine
emits an output determined by its current state \emph{and} the input symbol just
consumed. Both are automata reading a stream, and for both the only candidate cut is
between the machine and the stream --- again heterogeneous, again no well-sorted
$K(P,E)$. On the first axis, then, the three machines fail identically, and one might
think a single non-example would have sufficed.

It does not, because the pair discriminates along a \emph{second} axis: the migration
classification of Remark~\ref{ob:rem:migration} below. Were one to force an interaction-cut
reading on these machines, the Moore machine's output is a function of the state alone,
so nothing of the environment reaches the emitted value; the contraction would be pure
release, the null-migration corner occupied by CCS. The Mealy machine's output depends
on the consumed symbol, so the environment's datum does reach the residual; the
contraction would migrate. Two machines that fail interactivity for the same reason
nevertheless sit at different points of the classification that organises the successes.

\begin{nonexample}[Moore and Mealy machines]
Neither transducer has a naive presentation as an interactive GSLT, for the reason of
Non-example~\ref{ob:nonex:turing}. Under a forced reading they would occupy the
null-migration and binding-migration positions respectively.
\end{nonexample}

\section{Presentation versus encoding}

The reader will object that Turing machines can obviously be \emph{run} inside a process
calculus, and that the objection above is therefore pedantic. The objection is correct
and the distinction it presses on is the important one.

There is a difference between a theory \emph{having} a presentation and a theory
\emph{admitting an encoding into} one that does. Encoding a Turing machine into the rho
calculus --- reifying the tape as a collection of processes communicating on channels,
and the head as a process that moves among them --- produces a rho term whose behaviour
simulates the machine. That is a genuine and useful construction. But the result is a
statement about rho, not about
Turing machines. The interaction sites in the encoded term are rho's sites, and the
questions one can ask of them are questions about rho contexts.

This distinction is the entire subject of Chapter~\ref{ch:catgslt}, where the
morphisms of the category $\catGSLT$ turn out to be exactly the well-behaved encodings, and
where the conditions under which an encoding tells you something about the source rather
than merely about the target are isolated.

%===============================================================================
